\noindent Interpreting genetic variants has long been a central problem of clinical genetics and biomedical research. A single base change in a gene such as \textit{BRCA1} can be harmless, or it can sharply raise the risk of hereditary breast and ovarian cancer. Expert curation under the ACMG/AMP guidelines is reliable but slow, and it is least effective for rare variants that have never been observed before. With advances in Artificial Intelligence (AI), genomic language models such as Evo~2 have emerged as powerful tools that can score a mutation without any disease-specific training, but they are distributed as research code that needs a data-centre GPU and programming skills.

This project presents Evo2Serve, a full-stack web application that makes Evo~2-7B usable from a browser. A user selects a human genome assembly, browses chromosomes or searches a gene, reads the reference sequence, lists the known ClinVar variants, and enters a substitution, deletion, insertion or deletion--insertion. The model is served from a serverless NVIDIA H100 GPU on Modal through an authenticated FastAPI endpoint, and the interface is built with Next.js, React and TypeScript. For each request the system extracts an 8\,kb window around the variant, scores the reference and mutated sequences, and converts the score difference into a ``likely pathogenic'' or ``likely benign'' label with a confidence value, using a threshold calibrated on \textit{BRCA1} saturation genome editing data. The prediction is shown next to the expert ClinVar classification. On the 500 calibration variants the zero-shot model reached an AUROC of approximately 0.73, against 0.878 for Align-GVGD, 0.791 for CADD and 0.656 for SIFT computed on the same rows.

Because a prediction is only as good as the coordinate and the sequence it is computed on, the project also verifies the deployed pipeline: the composition of the calibration data, the context window, the genome assembly handling, and the two request validators over 55 generated cases. The verification identifies eight concrete improvements, the most important being to carry the reference allele in every request and to recalibrate the threshold on the full dataset. Evo2Serve is intended for education and research triage; its labels are computational evidence and not a clinical diagnosis.

\noindent\textbf{Keywords}\ \ - Genetic Variant Analysis, Variant Effect Prediction, Genomic Language Model, Evo 2, BRCA1, ClinVar, Saturation Genome Editing, ACMG/AMP Guidelines, Serverless GPU, FastAPI, Next.js, Web Application.
